\section{Introduction}
\label{sec:introduction}

The little finitistic dimension of a finite-dimensional algebra $A$ is the
supremum of the projective dimensions of the finitely generated $A$-modules of
finite projective dimension. Bass asked whether it is always
finite~\cite{Bas60}; see~\cite{ZHui95} for the history of the resulting
finitistic dimension conjecture. In September 2026, OpenAI released preprints
that claim counterexamples. The first~\cite{OAI26findim} constructs a
finite-dimensional complex algebra of infinite little finitistic dimension. The
second~\cite{OAI26ar} constructs a counterexample to the Auslander--Reiten
conjecture over the field $\mathbb{F}_2(q,H_1,H_2)$; by a theorem of Auslander
and Reiten~\cite[Theorem~1.1(b)]{AR75a}, together with
\Cref{thm:strong-nakayama}, such a counterexample yields a further
algebra of infinite little finitistic dimension.

This article reconstructs both constructions with proofs written out in full
(not verified by a human), organised
around general statements, some more general than the applications require. It
also gives criteria and obstructions that explain the choices in the
constructions. All proofs were written and checked by AI models, as
described in \Cref{app:verification,app:ai-declaration}; no human has checked
them, and the article has not been peer reviewed. The repository that
accompanies the preprints poses the existence statements of both constructions
as challenges in the proof assistant Lean and provides formal proofs of them;
\Cref{app:verification} compares these statements with the present article and
records what has been checked.

\subsection*{The main construction}

A \emph{selection functor}
${H=\Psi\otimes_R-}$ on finite-dimensional modules over an algebra $R$, possibly
infinite-dimensional, annihilates modules at arbitrarily late times
(\Cref{def:selection-data}). A bounded complex $P$ of bimodules over a
finite-dimensional algebra $B$ of global dimension at most two realises $H$, in
the
sense that ${P\Lotimes[B]M(Y)\cong M(HY)\oplus M(HY)[3]}$ for a fully faithful
embedding $M$ (\Cref{thm:realisation}). The complex $P$ is simulated by an
ordinary bimodule $X$ over an algebra $\Delta$ of finite global dimension, and
the trivial extension ${\Delta\ltimes X}$ turns the extinction times of $H$ into
projective dimensions (\Cref{prop:simulation,coro:extinction}). The result is
the following theorem, proved in \Cref{sec:simulation}; the main preprint
treats the case of a central idempotent and an endomorphism of~$R$.

\begin{theorem*}[\Cref{thm:main}]
  Let $(R,\Psi)$ be selection data over a field, with $R$ finitely presented
  and with unbounded extinction. Then there is a finite-dimensional algebra
  ${A=\Delta\ltimes X}$, with $\Delta$ of finite global dimension, such that
  ${\operatorname{findim}A=\infty}$.
\end{theorem*}

Selection data with unbounded extinction exist: the complex group algebra of a
finitely presented group with a family of independent central involutions,
shifted by an automorphism, provides them (\Cref{thm:selection-unbounded}),
which gives the counterexample of the main preprint (\Cref{coro:main}). The
selection data must be infinite-dimensional in an essential way: their rank
functions on finite-dimensional modules must span an infinite-dimensional space
(\Cref{prop:rank-obstruction}), which excludes finite-dimensional algebras,
commutative noetherian algebras and universal localisations of
finite-dimensional hereditary algebras
(\Cref{coro:rank-obstruction}). The realisation is the only step that is not
explicit: it lifts a diagram from a Verdier quotient by the calculus of
fractions, and the argument gives no bound on the length $l$ of the complex $P$,
hence on the number ${3(l+2)}$ of simple modules of the final algebra
(\Cref{rem:uncontrolled-length,rem:what-is-explicit}). The argument works with
an object ${U\oplus U[3]}$ instead of the summand $U$ that encodes $\Psi$: the
summand $U$ need not lie in the Verdier quotient, while ${U\oplus U[3]}$ does
(\Cref{prop:odd-double}), and with the shift $3$ the obstruction to splitting
the realisation lies in an $\operatorname{Ext}^4$ over an algebra of global
dimension at most two (\Cref{rem:why-three}).

\subsection*{The counterexample to the Auslander--Reiten conjecture}

The second construction uses the following result over an arbitrary field
(\Cref{thm:conversion}): for a symmetric algebra $E$, a
bimodule $F$ projective on both sides, a module $S$ and a stable morphism
${v\colon S\to FS}$, if the maps ${(g,j)\mapsto F(g)v-v[a]j}$ between Tate
cohomology groups are bijective in positive degrees and surjective with nonzero
kernel in degree zero, then a module over the triangular algebra
$\left(\begin{smallmatrix}E&0\\F&E\end{smallmatrix}\right)$ is a counterexample
to the Auslander--Reiten conjecture. The construction uses a ten-dimensional
algebra $C$, whose simple module has an aperiodic resolution, the symmetric
algebra ${T=C\ltimes DC}$, over which this simple module has polynomial Ext
algebra on a generator of degree three, and the tensor square ${E=T\otimes T}$,
over which two cones and two twists produce the bimodule $F$
(\Cref{sec:ar-route}). With \Cref{thm:ar-to-findim} this gives an algebra of
infinite little finitistic dimension with nine simple modules
(\Cref{coro:ar-findim}); the triangular algebra used to obtain it has eight
simple modules and dimension of the order of $10^{35}$ (\Cref{rem:ar-sizes}).

\subsection*{Criteria and obstructions}

The work that led to this article began with a search for a counterexample that
could be verified by hand. The search did not succeed, and its results are the
general statements of \Cref{sec:criteria,sec:extinction,sec:obstructions}.
Failure of the strong Nakayama conjecture produces explicit modules of every
finite projective dimension (\Cref{thm:strong-nakayama}); simple modules
witnessing this failure always come from a corner of the algebra and a
counterexample to the Auslander--Reiten conjecture over that corner
(\Cref{prop:simple-witness});
and finite projective dimensions, as well as finite extinction times, are
bounded on modules of bounded dimension
(\Cref{prop:bounded-dimension,prop:bounded-extinction}), where the assertion
about extinction times is for trivial extensions of algebras of finite global
dimension. For the construction of the counterexample to the Auslander--Reiten
conjecture, the analogous construction over $T$
itself, with one cone instead of two, fails in degree zero
whenever the simple module has polynomial Ext algebra on a generator of degree
three (\Cref{prop:one-factor,coro:one-factor}); the proof uses Tate duality,
which also forces the vanishing of the corresponding Toda brackets
(\Cref{thm:tate-obstruction}). Whether a counterexample with a small explicit
presentation exists remains open.

\subsection*{Structure}

\Cref{sec:preliminaries} fixes conventions and recalls Tate duality and Toda
brackets. \Cref{sec:criteria} contains the general criteria. The main
construction occupies \Cref{sec:extinction,sec:selection,sec:realisation,%
sec:simulation}: trivial extensions and extinction, selection data and the
group, the realisation by a bimodule complex, and the simulation with the main
theorem. The counterexample to the Auslander--Reiten conjecture is constructed
in \Cref{sec:conversion,sec:ar-route},
and \Cref{sec:obstructions} discusses obstructions to smaller constructions.
\Cref{app:computations} describes the computations used,
\Cref{app:verification} the verification and formalisation record,
\Cref{app:companion} the third preprint of the release, and
\Cref{app:ai-declaration} the use of AI in the preparation of this article.
